% Based on templates/lecture-example.tex.
% Related lecture: SC2005-L11, 2026-09-17.
% Sources: Part5 (Deadlocks).pdf, slides 5.4, 5.10--5.12, 5.22.
% E01--E04: source examples; E05: source diagram with explicitly added proof.
% No conversational concept-check questions are included.
% Human verified: lecture discussion completed; source graphs visually checked.

\exercisemeta
  {SC2005-L11-EX}
  {2026-09-17}
  {Part 5，pp.~5.4、5.10--5.12、5.22（PDF 页码相同）}
  {讲义例题的资源实例、箭头方向与完成顺序已核对；排序证明为补充推导}
  {原始讲义与解答已核对（2026-09-17）}

\exercisequestion{SC2005-L11-E01}{两个 Semaphore 的危险交错}

\textbf{来源：}讲义 p.~5.4，Example 2。\quad
\textbf{对应：}\relatedtopic{SC2005-L11-T01}{Deadlock 与资源分配模型}。

设 $A=B=1$。$P_0$ 依次执行 \texttt{wait(A)}、\texttt{wait(B)}；
$P_1$ 则依次执行 \texttt{wait(B)}、\texttt{wait(A)}。两者都在拿齐资源并完成操作后才释放。解释 context switch 如何导致死锁。

\begin{checkedsolution}
一种危险的交错为
\[
\begin{aligned}
  &P_0:\texttt{wait(A)} &&\text{成功，持有 }A;\\
  &P_1:\texttt{wait(B)} &&\text{成功，持有 }B;\\
  &P_1:\texttt{wait(A)} &&\text{阻塞，等待 }P_0;\\
  &P_0:\texttt{wait(B)} &&\text{阻塞，等待 }P_1.
\end{aligned}
\]
两者都到不了后面的释放操作。若某个进程先拿齐、完成并释放，则另一种交错可能顺利结束；一次成功运行不能证明该获取顺序安全。讲义 Example 1 的两种 tape drives 也是相同的“各持一种、再等另一种”结构。
\end{checkedsolution}

\exercisequestion{SC2005-L11-E02}{资源分配图：有等待但没有死锁}

\textbf{来源：}讲义 p.~5.10。\quad
\textbf{对应：}\relatedtopic{SC2005-L11-T01}{Deadlock 与资源分配模型}。

下图重绘了讲义的三个状态，圆为 process，矩形内每个点为一个 resource instance。红色箭头是 request，蓝色箭头是 assignment；绿色标出本例可先完成的 processes。图 (a) 中能否满足所有请求？给出一个完成顺序。

\begin{center}
  \resizebox{0.99\textwidth}{!}{% Original TikZ rendering of the graph states in Part 5, slides 5.10--5.12.
% Retains all resource counts and directed edges, including unused R4.
\begin{tikzpicture}[
  >=Latex,
  proc/.style={circle,draw=noteBlue,thick,fill=blue!5,minimum size=7mm,inner sep=1pt,font=\small},
  ready/.style={proc,draw=green!45!black,fill=green!8},
  res/.style={rectangle,draw=black!70,thick,fill=black!4,minimum width=7mm,minimum height=7mm},
  request/.style={-{Latex[length=2mm]},thick,draw=ntuRed},
  assigned/.style={-{Latex[length=2mm]},thick,draw=noteBlue},
  every label/.style={font=\small}
]
  \foreach \panel/\shift in {a/0,b/5.1} {
    \begin{scope}[xshift=\shift cm]
      \node[font=\small\bfseries] at (1.8,4.85)
        {\if a\panel (a) 无环，无死锁\else (b) 有环，有死锁\fi};
      \node[proc] (p1\panel) at (0.15,2.55) {$P_1$};
      \node[proc] (p2\panel) at (1.85,2.55) {$P_2$};
      \if a\panel
        \node[ready] (p3\panel) at (3.55,2.55) {$P_3$};
      \else
        \node[proc] (p3\panel) at (3.55,2.55) {$P_3$};
      \fi
      \node[res,label=above:$R_1$] (r1\panel) at (0.95,3.95) {};
      \node[res,label=above:$R_3$] (r3\panel) at (2.70,3.95) {};
      \node[res,minimum height=10mm,label=below:$R_2$] (r2\panel) at (0.95,0.95) {};
      \node[res,minimum height=13mm,label=below:$R_4$] (r4\panel) at (2.70,0.80) {};
      \coordinate (r2upper\panel) at (0.95,1.15);
      \coordinate (r2lower\panel) at (0.95,0.75);
      \fill (r1\panel) circle (1.4pt);
      \fill (r3\panel) circle (1.4pt);
      \fill (r2upper\panel) circle (1.4pt);
      \fill (r2lower\panel) circle (1.4pt);
      \foreach \y in {0.45,0.80,1.15} \fill (2.70,\y) circle (1.4pt);
      \draw[request] (p1\panel.north east) -- (r1\panel.south west);
      \draw[assigned] (r1\panel.center) -- (p2\panel.north west);
      \draw[request] (p2\panel.north east) -- (r3\panel.south west);
      \draw[assigned] (r3\panel.center) -- (p3\panel.north west);
      \draw[assigned] (r2lower\panel) -- (p1\panel.south east);
      \draw[assigned] (r2upper\panel) -- (p2\panel.south west);
      \if b\panel
        \draw[request] (p3\panel.south west) -- (r2\panel.east);
      \fi
      \node[font=\small] at (1.80,-0.45)
        {\if a\panel $P_3\to P_2\to P_1$\else 新增 $P_3\to R_2$\fi};
    \end{scope}
  }
  \begin{scope}[xshift=10.4cm]
    \node[font=\small\bfseries] at (1.9,4.85) {(c) 有环，无死锁};
    \node[proc] (p1c) at (0.1,2.45) {$P_1$};
    \node[ready] (p2c) at (3.5,4.05) {$P_2$};
    \node[proc] (p3c) at (3.5,2.45) {$P_3$};
    \node[ready] (p4c) at (3.5,0.60) {$P_4$};
    \node[res,minimum height=10mm,label=above:$R_1$] (r1c) at (1.45,3.70) {};
    \node[res,minimum height=10mm,label=below:$R_2$] (r2c) at (1.45,0.95) {};
    \coordinate (r1top) at (1.45,3.90);
    \coordinate (r1bottom) at (1.45,3.50);
    \coordinate (r2top) at (1.45,1.15);
    \coordinate (r2bottom) at (1.45,0.75);
    \foreach \point in {r1top,r1bottom,r2top,r2bottom} \fill (\point) circle (1.4pt);
    \draw[request] (p1c.north east) -- (r1c.south west);
    \draw[assigned] (r1top) -- (p2c.south west);
    \draw[assigned] (r1bottom) -- (p3c.north west);
    \draw[request] (p3c.south west) -- (r2c.north east);
    \draw[assigned] (r2top) -- (p1c.south east);
    \draw[assigned] (r2bottom) -- (p4c.north west);
    \node[font=\small] at (1.9,-0.45) {$P_2\to P_4\to P_3\to P_1$};
  \end{scope}
\end{tikzpicture}
}
  \par\small 根据讲义 pp.~5.10--5.12 的节点、实例数和边自行重绘；不增加额外资源需求。
\end{center}

\begin{checkedsolution}
图 (a) 中 $P_1$ 持有一个 $R_2$、等待 $R_1$；
$P_2$ 持有 $R_1$ 和另一个 $R_2$、等待 $R_3$；
$P_3$ 持有 $R_3$、没有新的请求。故可按
\[
  P_3\longrightarrow P_2\longrightarrow P_1
\]
完成：$P_3$ 先释放 $R_3$，$P_2$ 随后完成并释放 $R_1$，最后满足 $P_1$。这里假定图中列出的需求被满足后，进程能完成并归还资源；等待并不等于死锁。
\end{checkedsolution}

\exercisequestion{SC2005-L11-E03}{资源分配图：新增请求形成死锁}

\textbf{来源：}讲义 p.~5.11。\quad
\textbf{对应：}\relatedtopic{SC2005-L11-T01}{Deadlock 与资源分配模型}。

在图 (a) 的分配状态上增加 $P_3\rightarrow R_2$，得到图 (b)。判断是否死锁，并指出相应的环。

\begin{checkedsolution}
$R_2$ 的两个实例分别由 $P_1,P_2$ 持有，而两者又分别等待 $R_1,R_3$。原本可先完成的 $P_3$ 也在等待，相关进程都无法释放资源。可见两个环：
\[
\begin{aligned}
  &P_1\rightarrow R_1\rightarrow P_2\rightarrow R_3
    \rightarrow P_3\rightarrow R_2\rightarrow P_1,\\
  &P_2\rightarrow R_3\rightarrow P_3\rightarrow R_2\rightarrow P_2.
\end{aligned}
\]
系统存在死锁。虽然 $R_4$ 的三个实例均空闲，但它是不同类型，不能代替被请求的资源。$R_2$ 有多个实例也没有自动解决问题，因为其全部实例都被阻塞链中的进程占着。
\end{checkedsolution}

\exercisequestion{SC2005-L11-E04}{资源分配图：有环但没有死锁}

\textbf{来源：}讲义 p.~5.12。\quad
\textbf{对应：}\relatedtopic{SC2005-L11-T01}{Deadlock 与资源分配模型}。

图 (c) 中 $R_1$ 有两个实例，分给 $P_2,P_3$；$R_2$ 有两个实例，分给 $P_1,P_4$。
$P_1$ 请求 $R_1$，$P_3$ 请求 $R_2$。解释为何存在环，却仍可让所有进程完成。

\begin{checkedsolution}
图中确有
\[
  P_1\rightarrow R_1\rightarrow P_3\rightarrow R_2\rightarrow P_1,
\]
但环外的 $P_2,P_4$ 没有等待请求，可以分别释放 $R_1,R_2$ 的另一实例。讲义给出的一个完成顺序是
\[
  P_2\longrightarrow P_4\longrightarrow P_3\longrightarrow P_1.
\]
前两者完成后，$P_3$ 可获得一个 $R_2$ 并完成，随后 $P_1$ 也能完成。该顺序并不唯一：$P_2$ 释放 $R_1$ 后，$P_1$ 也已可能继续。

关键不是“多实例使图上的环消失”，而是 \textbf{同类的替代实例可解除某个请求}。只有当它空闲，或其 holder 能先完成并释放时才有帮助。
\end{checkedsolution}

\exercisequestion{SC2005-L11-E05}{筷子排序为何排除循环等待}

\textbf{来源：}讲义 p.~5.22，Chopstick Ordering 图示；以下反证为补充推导。\quad
\textbf{对应：}\relatedtopic{SC2005-L11-T02}{Deadlock Conditions 与 Prevention}。

五根筷子编号为 $0,\ldots,4$。$A,B,C,D$ 分别需要相邻的
$(0,1),(1,2),(2,3),(3,4)$，$E$ 需要 $(4,0)$。所有人均按较小编号在先获取。说明各人的顺序及其防止 circular wait 的理由。

\begin{center}
  \resizebox{0.87\textwidth}{!}{% Original ordering diagram based on Part 5, slide 5.22.
% Arrows encode each philosopher's acquisition order, NOT RAG edges.
\begin{tikzpicture}[
  >=Latex,
  resource/.style={draw=noteBlue,thick,fill=blue!5,rounded corners=2pt,minimum size=9mm,font=\large},
  order/.style={-{Latex[length=2.6mm]},very thick,draw=noteBlue},
  every node/.append style={font=\small}
]
  \foreach \i in {0,...,4} \node[resource] (c\i) at (2.6*\i,0) {$\i$};
  \foreach \i/\j/\name in {0/1/A,1/2/B,2/3/C,3/4/D}
    \draw[order] (c\i) -- node[above] {$\name$} (c\j);
  \draw[order,draw=ntuRed] (c0.south) to[out=-60,in=-120]
    node[below,text=ntuRed] {$E:\ 0\rightarrow4$，不能回绕成 $4\rightarrow0$} (c4.south);
\end{tikzpicture}
}
  \par\small 根据讲义图示整理为全局编号顺序；这里的箭头表示同一人的先后获取顺序，不是 RAG 的分配边。
\end{center}

\begin{checkedsolution}
顺序为 $A:0\rightarrow1$、$B:1\rightarrow2$、$C:2\rightarrow3$、
$D:3\rightarrow4$、$E:0\rightarrow4$。特别地，$E$ 不能按圆桌方向取成 $4\rightarrow0$。

假设仍存在等待环。沿环依次记录进程持有、并使前一进程等待的资源编号 $r_0,\ldots,r_k$。每位持有者只能继续申请更大编号，因而
\[
  r_0<r_1<\cdots<r_k<r_0,
\]
矛盾。所以统一 acquisition order（获取顺序）排除了 circular wait。证明不要求统一释放顺序，也不保证 starvation freedom；所有相关参与者必须遵守同一获取规则。
\end{checkedsolution}
